\expandafter\def\csname HistoricalNoteText14\endcsname{Es ist leicht zu zeigen, dass die gleiche Reihenfolge auch für den Wert $r=R$ selbst noch gilt, nur mit der Einschränkung, dass dann in ganz besonderen Fällen ein Übergangswert von $\rho$ (also einer, bei dem $T$ verschwindet oder $U=0$ wird) mit einer der Größen $-\omega,\ \omega,\ 3\omega,\ 5\omega$ usw. zusammenfallen kann. Für alle größeren Werte von $r$ muss dagegen jeder Übergangswert von $\rho$ \emph{zwischen} zwei dieser Größen liegen. Ich halte mich jedoch damit nicht auf, da es für unseren Zweck ausreicht, nachgewiesen zu haben, dass jene Reihenfolge von irgendeinem Wert von $r$ an gilt.}
